\documentclass[11pt,a4paper]{article}
\usepackage[utf8]{inputenc}
\usepackage[T1]{fontenc}
\usepackage[spanish]{babel}
\usepackage{lmodern}
\usepackage{geometry}
\geometry{margin=2.2cm}
\usepackage{xcolor,listings,booktabs,hyperref,longtable}
\usepackage{fancyhdr}
\hypersetup{colorlinks=true,linkcolor=blue,urlcolor=blue}
\lstset{basicstyle=\ttfamily\small,breaklines=true,frame=single,columns=fullflexible,showstringspaces=false,inputencoding=utf8,literate={á}{{\'a}}1 {é}{{\'e}}1 {í}{{\'i}}1 {ó}{{\'o}}1 {ú}{{\'u}}1 {Á}{{\'A}}1 {É}{{\'E}}1 {Í}{{\'I}}1 {Ó}{{\'O}}1 {Ú}{{\'U}}1 {ñ}{{\~n}}1 {Ñ}{{\~N}}1 {ü}{{\"u}}1 {Ü}{{\"U}}1}
\pagestyle{fancy}
\fancyhf{}
\lhead{IFCD107}
\rhead{Especialista en Inteligencia Artificial}
\cfoot{\thepage}
\setlength{\headheight}{14pt}
\title{API REST de regresión con Flask\\\large Arquitectura, endpoints, flujos y decisiones de diseño}
\author{IFCD107 -- Especialista en Inteligencia Artificial\\Proyecto práctico de API y aprendizaje automático}
\date{Octubre de 2026}
\begin{document}
\maketitle
\thispagestyle{fancy}
\begin{abstract}
Este proyecto documenta una API REST Flask para regresión: ajuste polinómico, búsqueda de hiperparámetros, KNN con dos esquemas de validación, Random Forest, SVR, AdaBoost y Gradient Boosting. La guía explica el contrato de cada endpoint, sus formatos, flujos de datos, artefactos, autenticación y las decisiones que organizan la implementación. La colección Postman contiene solicitudes de ejemplo.
\end{abstract}

\section{Descripción de la solución}
La aplicación reúne regresión polinómica sobre pares $(x,y)$ y modelos tabulares con variables numéricas y categóricas: KNN, Random Forest, SVR, AdaBoost y Gradient Boosting. La búsqueda polinómica compara configuraciones por validación cruzada. El ejemplo de dominio representa expedientes de terrazas con datos sintéticos; no contiene registros reales ni reproduce estadísticas o plazos del Ayuntamiento de Madrid. Las ejecuciones guardan los artefactos correspondientes bajo una referencia y un identificador único.

\begin{center}
\begin{tabular}{ll}\toprule
Componente & Responsabilidad\\\midrule
\texttt{app.py} & Endpoints HTTP, validación, carpetas, persistencia y respuesta JSON.\\
\texttt{regression/polynomial\_regressor.py} & Ajuste matemático, métricas, CSV de salida y gráfica.\\
\texttt{regression/random\_forest\_regressor.py} & Preprocesamiento mixto, entrenamiento RF e importancia de variables.\\
\texttt{scripts/generate\_terrace\_cases.py} & Generador determinista de las 1.000 filas sintéticas.\\
\texttt{ejemplos/datos\_ejemplo\_xy.csv} & Datos con columnas $x,y$ para los endpoints polinómicos.\\
\texttt{postman/} & Colección importable con health, autenticación Bearer, solicitudes JSON/CSV de ambos algoritmos.\\
\texttt{runs/} & Artefactos generados, separados por referencia y ejecución.\\\bottomrule
\end{tabular}
\end{center}

\section{Estructura de directorios}
\begin{lstlisting}
proyecto_regresion_polinomica/
  app.py
  api_functions.py
  regression/polynomial_regressor.py
  regression/random_forest_regressor.py
  ejemplos/datos_ejemplo_xy.csv
  ejemplos/datos_ejemplo_complejo.csv
  postman/Regresion_Polinomica.postman_collection.json
  documentacion/guia_proyecto.tex
  runs/<referencia>/
    graficas/<run_id>_ajuste.png
    csv_entrada/<run_id>.csv
    logs_salida/<run_id>.log
    csv_salida/<run_id>_predicciones.csv
    modelos/<run_id>_modelo.pkt
\end{lstlisting}
La referencia admite letras, números, punto, guion y guion bajo, con máximo de 64 caracteres. Se rechazan rutas o separadores para evitar escapar del directorio de ejecución. Cada petición genera un identificador único, llamado \texttt{run\_id}; así se guardan varias ejecuciones sin sobrescribir los archivos de la misma referencia.

\section{Instalación y arranque}
Se recomienda Python 3.10 o superior. Desde la carpeta raíz del proyecto:
\begin{lstlisting}
python -m venv .venv
\end{lstlisting}
Activar el entorno: en Windows PowerShell, \texttt{.venv\\Scripts\\Activate.ps1}; en macOS/Linux, \texttt{source .venv/bin/activate}. Después:
\begin{lstlisting}
python -m pip install -r requirements.txt
\end{lstlisting}
La API carga automáticamente las variables del archivo local \texttt{.env} al arrancar. En esta copia se incluye un \texttt{.env} local con una contraseña aleatoria para el usuario \texttt{admin} y un secreto de firma. Los valores no se reproducen en esta guía. Para otro entorno, copia \texttt{.env.example} a \texttt{.env} y configura credenciales propias. El archivo \texttt{.env} está excluido de Git y del ZIP distribuible. Copia el valor \texttt{API\_ADMIN\_PASSWORD} a la variable \texttt{password} de Postman. Mantén \texttt{API\_TOKEN\_SECRET} estable entre procesos y reinicios.

La API queda disponible localmente en \url{http://127.0.0.1:5000}. \url{http://127.0.0.1:5000/api/v1/health} es público y devuelve el estado del servicio.

\section{Flujo de ejecución}
Para seguir el programa desde el principio, piense en cada petición como una cadena de pasos: Postman envía datos, Flask valida la petición, una función de \texttt{api\_functions.py} ejecuta el modelo, y finalmente Flask devuelve JSON con el resultado y las rutas de los archivos creados.

\begin{enumerate}
\item \textbf{Arranque.} Al ejecutar \texttt{python app.py}, Flask crea la aplicación y registra sus rutas. La ruta pública \texttt{/api/v1/health} permite comprobar que el servidor responde.
\item \textbf{Autenticación.} El cliente envía usuario y contraseña a \texttt{/api/v1/auth/token}. Si coinciden con la configuración, la API devuelve un token firmado. En las peticiones de entrenamiento, el decorador \texttt{@require\_token} comprueba el token antes de ejecutar la ruta.
\item \textbf{Lectura y validación.} Las rutas polinómicas aceptan puntos JSON o archivos CSV con encabezados \texttt{x,y}. Las rutas Random Forest reciben JSON para el conjunto de ejemplo o CSV multipart. Flask valida referencia, formato, tamaño y parámetros antes de iniciar el entrenamiento.
\item \textbf{Entrenamiento.} La ruta llama a una función en \texttt{api\_functions.py}. Esa función divide los datos, construye el estimador o pipeline y calcula predicciones y métricas.
\item \textbf{Archivos y respuesta.} La función crea un \texttt{run\_id}, guarda los artefactos bajo \texttt{runs/<referencia>/} y devuelve el identificador, métricas y rutas en una respuesta JSON.
\end{enumerate}

\subsection{Flujo de regresión polinómica}
\texttt{POST /api/v1/regresion/polinomica} recibe JSON con \texttt{degree} y una lista \texttt{data}. La variante \texttt{/csv} recibe \texttt{reference}, \texttt{degree} y un archivo CSV. Ambas terminan en \texttt{run\_model}: esta función separa entrenamiento y prueba, crea \texttt{PolynomialRegressor}, ajusta los coeficientes con NumPy y calcula predicciones y métricas para la partición de prueba. Después guarda el modelo, las predicciones, la gráfica y el log.

\subsection{Flujo de RandomizedSearchCV}
\texttt{POST /api/v1/regresion/polinomica/randomizedsearch} admite los dos formatos: JSON con \texttt{data}, o \texttt{multipart/form-data} con \texttt{file} (CSV con columnas \texttt{x,y}). En ambos casos también se pueden enviar \texttt{n\_iter}, \texttt{cv} y \texttt{random\_state}. El campo opcional \texttt{w} es una lista de pesos, uno por cada fila.

La ruta valida y convierte la entrada, y llama a \texttt{run\_polynomial\_search}. Esta función reserva primero el 20\% de las filas como prueba final. En el 80\% restante, \texttt{RandomizedSearchCV} muestrea combinaciones del diccionario \texttt{POLYNOMIAL\_PARAM\_DISTRIBUTIONS}; para cada combinación, valida el resultado en \texttt{cv} pliegues. El \texttt{Pipeline} conecta el adaptador \texttt{PolynomialSearchEstimator} con scikit-learn, mientras que el adaptador llama a \texttt{PolynomialRegressor}. Cuando termina la búsqueda, la mejor configuración se vuelve a ajustar con todos los datos de entrenamiento y se evalúa una sola vez con el conjunto de prueba reservado.

La respuesta de búsqueda incluye \texttt{best\_parameters}, el RMSE promedio de validación cruzada, las métricas de prueba y las rutas de los archivos. El log también registra el diccionario de distribuciones y los parámetros del mejor estimador. \texttt{full} y \texttt{cov} controlan diagnósticos del ajuste; no son métricas de rendimiento, por lo que no cambian el criterio de selección (RMSE).

\subsection{Flujo de Random Forest}
Las rutas \texttt{/api/v1/random-forest/terrazas} y \texttt{/api/v1/random-forest/csv} llaman a \texttt{run\_random\_forest}. La clase \texttt{RandomForestRegressorModel} separa la columna objetivo de las predictoras, prepara las columnas numéricas y categóricas en un \texttt{ColumnTransformer}, y combina ese preprocesamiento con \texttt{RandomForestRegressor} en un \texttt{Pipeline}. El conjunto de prueba queda fuera del entrenamiento. La respuesta incluye métricas, las importancias principales y las rutas de los artefactos.

\subsection{Flujos KNN, SVR y boosting}
KNN recibe un CSV tabular y genera predicciones fuera de muestra con K-Fold o Leave-One-Out. En K-Fold, las filas se reparten en $k$ pliegues y cada pliegue se predice cuando está excluido del ajuste; Leave-One-Out repite el proceso dejando una sola fila fuera. Estas rutas calculan validación cruzada directamente. El pipeline KNN imputa y escala predictores antes de calcular distancias.

SVR, AdaBoost y Gradient Boosting comparten lectura CSV, el parámetro \texttt{target}, la separación train/test y el guardado. La función común \texttt{run\_tabular\_regressor} selecciona el constructor del algoritmo. SVR escala variables numéricas porque depende de distancias y márgenes; los modelos de árboles de boosting no requieren escalado. En todos los casos el preprocesamiento se ajusta únicamente con entrenamiento para evitar filtración desde test.

\section{Catálogo de endpoints}
Todas las rutas de entrenamiento usan \texttt{Authorization: Bearer <token>}. Las solicitudes CSV se envían como \texttt{multipart/form-data}; no se debe fijar manualmente el encabezado boundary. La ruta de salud y la emisión de token son públicas.

\begin{longtable}{p{0.23\linewidth}p{0.22\linewidth}p{0.44\linewidth}}
\toprule Método y ruta & Entrada & Función y resultado \\
\midrule
\endhead
GET \texttt{/api/v1/health} & Sin cuerpo & Comprueba disponibilidad; devuelve HTTP 200.\\
POST \texttt{/api/v1/auth/token} & JSON: \texttt{username}, \texttt{password} & Emite \texttt{access\_token}, tipo Bearer y expiración; credenciales incorrectas: 401.\\
POST \texttt{/api/v1/regresion/polinomica} & JSON: \texttt{reference}, \texttt{degree}, \texttt{data:[\{x,y\}]} & Ajusta polinomio y evalúa test; devuelve coeficientes, métricas y artefactos.\\
POST \texttt{/api/v1/regresion/polinomica/csv} & Multipart: \texttt{reference}, \texttt{degree}, \texttt{file} con columnas \texttt{x,y} & Mismo ajuste usando CSV.\\
POST \texttt{/api/v1/regresion/polinomica/randomizedsearch} & JSON con \texttt{data} o multipart con \texttt{file}; \texttt{n\_iter}, \texttt{cv}, \texttt{random\_state}, opcional \texttt{w} & Busca parámetros mediante validación cruzada en train y evalúa una vez en test.\\
POST \texttt{/api/v1/random-forest/terrazas} & JSON: referencia, objetivo y opciones RF & Entrena con CSV sintético incluido; devuelve métricas e importancias.\\
POST \texttt{/api/v1/random-forest/csv} & Multipart: referencia, \texttt{file}, objetivo y opciones RF & Entrena Random Forest con datos aportados.\\
POST \texttt{/api/v1/regresion/knn/kfold} & Multipart: referencia, CSV, objetivo, \texttt{n\_neighbors}, \texttt{cv}, semilla & KNN con predicción cruzada K-Fold.\\
POST \texttt{/api/v1/regresion/knn/leave-one-out} & Multipart: referencia, CSV, objetivo, \texttt{n\_neighbors}, semilla & KNN Leave-One-Out, una predicción fuera de muestra por fila.\\
POST \texttt{/api/v1/regresion/svm/csv} & Multipart: CSV, objetivo, referencia y opciones SVR & Opciones: \texttt{kernel}, \texttt{C}, \texttt{epsilon}, \texttt{gamma}, \texttt{degree}.\\
POST \texttt{/api/v1/regresion/boosting/adaboost/csv} & Multipart: CSV, objetivo, referencia y opciones & AdaBoostRegressor: \texttt{n\_estimators}, \texttt{max\_depth}, \texttt{learning\_rate}, \texttt{loss}.\\
POST \texttt{/api/v1/regresion/boosting/gradient/csv} & Multipart: CSV, objetivo, referencia y opciones & GradientBoostingRegressor: parámetros anteriores y \texttt{subsample}.\\
\bottomrule
\end{longtable}

Los modelos tabulares aceptan además \texttt{test\_size} (0.2 por defecto) y \texttt{random\_state} (42), salvo KNN, que calcula validación cruzada. SVR admite kernels \texttt{linear}, \texttt{poly}, \texttt{rbf} o \texttt{sigmoid}; \texttt{gamma} admite \texttt{scale}, \texttt{auto} o valor numérico. Los valores por defecto y límites se validan en \texttt{regression\_model\_options} de \texttt{api\_functions.py}.

\subsection{Contrato de respuesta y errores}
Una ejecución correcta devuelve JSON con \texttt{reference}, \texttt{run\_id}, métricas, detalles propios del modelo y \texttt{artifacts}. El modelo polinómico informa grado y coeficientes; la búsqueda informa su mejor configuración; Random Forest incluye importancias; KNN describe folds; los demás modelos informan opciones y métricas. Los errores contienen el campo \texttt{error}.

Estados utilizados: 200 para resultado correcto; 400 para campos ausentes, opciones no válidas, CSV inválido o referencia no permitida; 401 para credenciales o Bearer inválidos/expirados; 413 para cuerpo o CSV demasiado grande; 422 cuando una entrada válida sintácticamente no permite completar el ajuste; y 500 para fallos inesperados.

\section{Decisiones de diseño}
\begin{description}
\item[Separación por responsabilidades.] \texttt{app.py} define el contrato HTTP; \texttt{api\_functions.py} concentra validación y orquestación; \texttt{regression/} implementa los modelos. El flujo se puede seguir desde la ruta hasta el estimador sin duplicar manejo HTTP.
\item[JSON para puntos, CSV para tablas.] Los puntos polinómicos son compactos como JSON. Los modelos de varias columnas usan CSV multipart, que preserva encabezados y admite archivos de curso.
\item[Pipeline para datos tabulares.] Imputación, codificación y estimador se ajustan juntos. En divisiones train/test y folds, las transformaciones aprenden solo de las filas de entrenamiento.
\item[Preprocesamiento.] Variables numéricas: imputación por mediana; categóricas: imputación por moda y One-Hot Encoding que tolera categorías nuevas. El identificador de expediente se conserva para trazabilidad cuando procede, pero se excluye de predictores.
\item[Escalado según algoritmo.] SVR y KNN necesitan escalado por su dependencia de distancias. Random Forest y árboles de boosting usan particiones y no requieren estandarización.
\item[Evaluación.] Polinomio, Random Forest, SVR y boosting separan test. La búsqueda polinómica selecciona por RMSE en folds de train y reserva test para la evaluación final. KNN genera directamente predicciones fuera de fold o Leave-One-Out.
\item[Reproducibilidad.] Las rutas aceptan semilla; la referencia agrupa ejecuciones y el identificador \texttt{run\_id} evita sobrescrituras.
\item[Autenticación central.] Un decorador Bearer protege las rutas de entrenamiento sin duplicar la comprobación; health permanece pública.
\item[Auditoría.] Se registra método, ruta, estado, usuario, IP y huella truncada del token; no se registra el token completo. Credenciales y secreto se configuran en \texttt{.env}.
\item[Serialización.] Pickle se usa para la demostración local. Solo se deben cargar modelos generados por una fuente confiable, pues un pickle manipulado puede ejecutar código.
\end{description}

\section{Prueba desde Postman}
\begin{enumerate}
\item En Postman, importar \texttt{postman/Regresion\_Polinomica.postman\_collection.json}.
\item Mantener \texttt{baseUrl=http://127.0.0.1:5000} y usar la referencia \texttt{terrazas\_madrid\_demo}.
\item Configurar las variables de colección \texttt{username=admin} y \texttt{password} con la clave definida en \texttt{API\_ADMIN\_PASSWORD}.
\item Ejecutar \textbf{Health}; debe responder HTTP 200.
\item Ejecutar \textbf{Obtener token}. La colección guarda automáticamente el token Bearer recibido. Credenciales incorrectas producen HTTP 401.
\item Ejecutar \textbf{Ajustar desde JSON}. La colección ya incluye datos sintéticos y grado 2. Postman establece \texttt{Content-Type: application/json}.
\item Para los datos polinómicos CSV, ejecutar \textbf{Ajustar desde CSV} y seleccionar \texttt{ejemplos/datos\_ejemplo\_xy.csv}. Postman crea el formulario multipart; no fijar manualmente el encabezado boundary.
\item Ejecutar \textbf{Regresión polinómica - RandomizedSearchCV} para buscar parámetros con puntos JSON, o \textbf{RandomizedSearchCV - cargar CSV} para adjuntar \texttt{ejemplos/datos\_ejemplo\_xy.csv}. Ambas solicitudes llaman a \texttt{/api/v1/regresion/polinomica/randomizedsearch}.
\item Ejecutar \textbf{Random Forest - CSV de terrazas incluido} para entrenar sin adjuntar fichero, o \textbf{Random Forest - cargar CSV} y seleccionar \texttt{ejemplos/datos\_ejemplo\_complejo.csv}. El objetivo por defecto es \texttt{y}; se excluye \texttt{expediente\_id} y se usan como predictores el resto de columnas.
\item Leer \texttt{artifacts} en la respuesta JSON. Las rutas son relativas al directorio del proyecto y están agrupadas bajo \texttt{runs/<reference>/}.
\end{enumerate}

\subsection{Obtención del token}
\begin{lstlisting}
POST /api/v1/auth/token
Content-Type: application/json

{ "username": "admin", "password": "<clave configurada>" }
\end{lstlisting}
La respuesta incluye \texttt{access\_token}, \texttt{token\_type=Bearer} y \texttt{expires\_in}. La duración predeterminada es una hora y se puede cambiar con \texttt{API\_TOKEN\_TTL\_SECONDS}. Los endpoints de entrenamiento requieren el encabezado Bearer; sin token, con token inválido o expirado responden HTTP 401.

\subsection{Solicitud JSON}
\begin{lstlisting}
POST /api/v1/regresion/polinomica
Authorization: Bearer <token>
Content-Type: application/json

{
  "reference": "expediente_demo_001",
  "degree": 2,
  "data": [
    {"x": -1, "y": 2.2}, {"x": 0, "y": 1.0},
    {"x": 1, "y": 1.8}, {"x": 2, "y": 6.2}
  ]
}
\end{lstlisting}

\subsection{Solicitud CSV}
\begin{lstlisting}
POST /api/v1/regresion/polinomica/csv
Authorization: Bearer <token>
Body: form-data
  reference = terrazas_madrid_demo
  degree    = 2
  file      = ejemplos/datos_ejemplo_xy.csv (columnas x,y)
\end{lstlisting}

\subsection{Solicitud de búsqueda polinómica}
La ruta JSON y la ruta de archivo comparten el endpoint \texttt{/api/v1/regresion/polinomica/randomizedsearch}. La diferencia está en el formato de la petición:
\begin{lstlisting}
# JSON: enviar puntos directamente en el cuerpo
POST /api/v1/regresion/polinomica/randomizedsearch
Authorization: Bearer <token>
Content-Type: application/json

{
  "reference": "busqueda_demo",
  "n_iter": 20,
  "cv": 5,
  "random_state": 42,
  "data": [
    {"x": 0, "y": 1}, {"x": 1, "y": 2}, {"x": 2, "y": 5},
    {"x": 3, "y": 10}, {"x": 4, "y": 17}, {"x": 5, "y": 26},
    {"x": 6, "y": 37}, {"x": 7, "y": 50}, {"x": 8, "y": 65},
    {"x": 9, "y": 82}
  ]
}

# CSV: en Postman elegir Body > form-data y adjuntar el fichero
reference   = busqueda_csv_demo
n_iter      = 20
cv          = 5
random_state = 42
file        = ejemplos/datos_ejemplo_xy.csv
\end{lstlisting}
El CSV debe tener los encabezados \texttt{x,y}. Si se usa el campo opcional \texttt{w}, se envía una lista JSON de pesos con tantos valores como filas tenga \texttt{data}; por ejemplo, \texttt{[1,1,0.5,...]}. Para JSON es un array normal; en multipart se escribe como texto con formato JSON.

\subsection{Entrenamiento Random Forest}
La ruta de demostración lee el CSV local \texttt{ejemplos/datos\_ejemplo\_complejo.csv} (1.000 filas sintéticas) y la ruta multipart admite otro CSV. Ambas requieren Bearer y guardan CSV de entrada, predicciones, gráfica, modelo \texttt{.pkt} y log en el mismo árbol \texttt{runs/<referencia>/}.

\begin{lstlisting}
POST /api/v1/random-forest/terrazas
Authorization: Bearer <token>
Content-Type: application/json

{
  "reference": "terrazas_rf_demo",
  "target": "y",
  "n_estimators": 200,
  "max_depth": null,
  "test_size": 0.2,
  "random_state": 42
}
\end{lstlisting}
La ruta \texttt{/api/v1/random-forest/csv} acepta \texttt{multipart/form-data}: \texttt{reference}, \texttt{file} y, opcionalmente, \texttt{target}, \texttt{n\_estimators}, \texttt{max\_depth}, \texttt{test\_size} y \texttt{random\_state}. Se imputan valores ausentes por mediana en variables numéricas y por moda en categóricas; estas últimas se codifican con one-hot encoding. El identificador de expediente no se usa como predictor. La división aleatoria fija una semilla y separa train/test; las métricas se calculan únicamente en test. La respuesta incluye $R^2$, RMSE, MAE e importancia de las 15 variables transformadas más relevantes. Los datos del ejemplo son sintéticos y las métricas no describen el rendimiento real de la tramitación municipal.

\section{Respuesta y archivos}
Las respuestas exitosas incluyen \texttt{reference}, \texttt{run\_id}, métricas y \texttt{artifacts}; también llevan parámetros, coeficientes, predicciones o importancias según el estimador. Los artefactos habituales son CSV de entrada, CSV de predicciones, gráfica, modelo serializado y log de ejecución. La respuesta KNN informa el método de validación y las métricas de sus predicciones fuera de muestra; los modelos de árboles pueden incluir importancia de variables.

Cada ejecución crea sus archivos dentro de \texttt{runs/<reference>/} en subcarpetas \texttt{csv\_entrada}, \texttt{csv\_salida}, \texttt{graficas}, \texttt{modelos} y \texttt{logs\_salida}. Las rutas se devuelven en \texttt{artifacts}. El registro \texttt{runs/\_auditoria/api\_audit.jsonl} añade método, ruta, estado, usuario, IP y huella truncada del token para auditoría.

El formato \texttt{.pkt} de este ejemplo utiliza serialización Python pickle. Solo se deben cargar archivos generados localmente y de confianza: pickle puede ejecutar código al cargar un fichero manipulado. Para intercambio entre sistemas, conviene guardar los coeficientes y metadatos en JSON en vez de cargar pickles externos.

\section{Validaciones y límites del ejemplo}
La API limita el cuerpo HTTP a 5 MB, los archivos CSV a 4 MB y las entradas polinómicas a 100.000 filas. Valida referencia, formato, finitud y parámetros; la búsqueda polinómica limita \texttt{n\_iter} a 1--80, \texttt{cv} a 2--10 y exige datos suficientes después de separar test. Devuelve HTTP 400 para entradas no válidas, 401 para credenciales/token no válidos, 413 para contenido demasiado grande, 422 cuando un ajuste no se puede completar y 500 para fallos inesperados. Health y emisión de token son públicas; los entrenamientos requieren Bearer.

\section{Notas sobre interpretación}
Random Forest no extrapola tendencias fuera del rango de entrenamiento de forma lineal y sus importancias no implican causalidad. La evaluación usa una partición aleatoria; no modela autocorrelación geográfica ni temporal.

La regresión polinómica es sensible al grado, al rango y escala de $x$, a valores atípicos y a extrapolación. $R^2$ y errores se calculan sobre los mismos datos usados para el ajuste; por sí solos no estiman generalización. Para uso predictivo real, separar entrenamiento y evaluación, comparar grados con validación y revisar residuos y supuestos del dominio.

Este caso sirve para demostrar el flujo técnico, no para estimar plazos reales de expedientes ni evaluar el desempeño de unidades municipales. Para ello harían falta datos autorizados, definición institucional del inicio/fin de tramitación, tratamiento de expedientes abiertos/censurados, análisis de estacionalidad y validación temporal fuera de muestra. Se puede regenerar el fichero con \texttt{python scripts/generate\_terrace\_cases.py}; la semilla fija mantiene la demo reproducible.

\appendix
\section{Mapa del código fuente}
La documentación describe el contrato funcional; el código ejecutable del proyecto es la referencia final de los detalles de implementación. Los comentarios de los módulos están redactados para orientar a personas que empiezan en Python.

\begin{longtable}{p{0.36\linewidth}p{0.53\linewidth}}
\toprule Archivo & Contenido principal \\
\midrule
\endhead
\texttt{app.py} & Fábrica Flask, definición de rutas, emisión/verificación de tokens, límites HTTP y auditoría.\\
\texttt{api\_functions.py} & Validación de referencia y entradas, conversión de CSV/JSON, ejecución de modelos, métricas y escritura de artefactos.\\
\texttt{regression/polynomial\_regressor.py} & Regresión polinómica y adaptador de búsqueda para scikit-learn.\\
\texttt{regression/knn\_regressor.py} & Pipeline KNN y preparación de la validación cruzada.\\
\texttt{regression/random\_forest\_regressor.py} & Pipeline Random Forest, predicciones e importancias.\\
\texttt{regression/tabular\_preprocessing.py} & Imputación, codificación categórica y escalado numérico compartido.\\
\texttt{regression/svm\_regressor.py} & Constructor del pipeline SVR.\\
\texttt{regression/adaboost\_regressor.py} & Constructor del pipeline AdaBoostRegressor.\\
\texttt{regression/gradient\_boosting\_regressor.py} & Constructor del pipeline GradientBoostingRegressor.\\
\texttt{postman/Regresion\_Polinomica.postman\_collection.json} & Peticiones reproducibles para health, token y endpoints de entrenamiento.\\
\bottomrule
\end{longtable}

\section{Flujo para probar la API}
\begin{enumerate}
\item Instalar dependencias desde \texttt{requirements.txt} y arrancar \texttt{python app.py}.
\item Importar la colección Postman y configurar \texttt{baseUrl}, \texttt{username} y la contraseña local.
\item Ejecutar Health y después Obtener token; la colección guarda el token para las rutas protegidas.
\item Elegir una petición según el formato: JSON para puntos polinómicos o multipart con archivo para datos tabulares.
\item Revisar el código HTTP y el cuerpo JSON, y localizar los archivos indicados en \texttt{artifacts} bajo \texttt{runs/<reference>/}.
\end{enumerate}

\end{document}
